\section{Results}
\label{sec:results}

\subsection{Central case: the trajectory the reference data imply}
\label{sec:results:base}

Under the reference values of every factor, the Peruvian system roughly quadruples
its installed generation capacity between 2025 and 2050, from \SI{14.06}{\giga\watt}
to \SI{63.85}{\giga\watt}, while the transmission network grows from
\SI{16180}{\kilo\meter} to \SI{46341}{\kilo\meter} and inter-zonal transfer
capability from \SI{1.60}{\giga\watt} to \SI{4.56}{\giga\watt}. The country average
electricity price falls from roughly \SI{210}{\USD\per\mega\watt\hour} in the late
2020s to \SI{124.5}{\USD\per\mega\watt\hour} by 2050.

The composition of that growth is the first result worth noting. Almost all of it is
onshore wind and solar: wind rises from \SI{1.02}{\giga\watt} to
\SI{34.83}{\giga\watt} and solar PV from \SI{0.48}{\giga\watt} to
\SI{8.33}{\giga\watt}, while thermal capacity grows only from \SI{7.25}{\giga\watt} to
\SI{9.65}{\giga\watt}, so non-thermal technologies hold 84.9\,\% of capacity by 2050
against 48.4\,\% at the start. The model reaches this without any renewable target,
mandate or carbon price: expansion is driven by the profitability rule alone. Thermal
plant is not retired --- the model has no endogenous retirement flow --- but it stops
being the marginal investment.

The price path is not monotone: it rises through the late 2020s, when demand grows
against a capacity base that cannot yet respond, and falls thereafter as wind and solar
enter. That is the classic capacity-cycle signature system dynamics models of
liberalised markets reproduce \citep{ford2001,arango2007,simshauser2010}, and it
appears here without being imposed.

\subsection{Which uncertainty governs the system}
\label{sec:results:ranking}

Figure~\ref{fig:tornado} reports the one-at-a-time sensitivity of each indicator. The
result is unambiguous: demand growth dominates by roughly an order of magnitude over
the next-ranked factor, and fuel prices --- the risk most often emphasised for a
gas-dependent system --- are the weakest driver of three of the four indicators.

\begin{figure*}[tb]
  \centering
  \includegraphics[width=\textwidth]{03_tornado.png}
  \caption{Tornado diagram. For each indicator, the bars show the change from the
  central case when one factor is set to the low and the high end of its range with
  the other three held at their reference value. Factors are ordered by swing,
  $|\text{high}-\text{low}|$. Values at January 2050; the price is the mean of the
  final twelve months.}
  \label{fig:tornado}
\end{figure*}

Table~\ref{tab:tornado} gives the numbers. A $\pm$20\,\% band on the 2050 demand level
moves required transmission from $-18.9$\,\% to $+23.8$\,\% and installed capacity from
$-24.0$\,\% to $+29.3$\,\%; the swing attributable to demand is nine times that of
generation lead times for transmission length. Fuel prices move the average price by
0.4\,\% in both directions --- not a sign error but a genuine non-monotonicity.

\begin{table}[tb]
  \centering
  \caption{One-at-a-time sensitivity. Change from the central case at January 2050
  with one factor at each end of its range. Factors ordered by swing within each
  indicator.}
  \label{tab:tornado}
  \scriptsize
  \begin{tabular}{@{}llrrr@{}}
    \toprule
    Indicator & Factor & Low & High & Swing \\
    \midrule
    \multirow{4}{*}{\shortstack[l]{Transmission\\ \si{\kilo\meter}, base \num{46341}}}
      & Demand growth        & $-18.9\%$ & $+23.8\%$ & \num{19794} \\
      & Generation lead time & $+6.4\%$  & $+1.7\%$  & \num{2210}  \\
      & Fuel price           & $+0.2\%$  & $+1.8\%$  & \num{729}   \\
      & Transmission lead time & $+1.8\%$ & $+0.7\%$ & \num{487}   \\
    \midrule
    \multirow{4}{*}{\shortstack[l]{Installed capacity\\ \si{\giga\watt}, base \num{63.85}}}
      & Demand growth        & $-24.0\%$ & $+29.3\%$ & \num{34.08} \\
      & Generation lead time & $+10.6\%$ & $+1.1\%$  & \num{6.09}  \\
      & Fuel price           & $+0.4\%$  & $+1.4\%$  & \num{0.64}  \\
      & Transmission lead time & $+1.4\%$ & $+1.2\%$ & \num{0.08}  \\
    \midrule
    \multirow{4}{*}{\shortstack[l]{Transfer capability\\ \si{\giga\watt}, base \num{4.56}}}
      & Demand growth        & $-8.4\%$  & $+15.4\%$ & \num{1.08}  \\
      & Transmission lead time & $+5.2\%$ & $-4.9\%$ & \num{0.46}  \\
      & Fuel price           & $+4.9\%$  & $+0.6\%$  & \num{0.19}  \\
      & Generation lead time & $-2.1\%$  & $-4.3\%$  & \num{0.10}  \\
    \midrule
    \multirow{4}{*}{\shortstack[l]{Average price\\ \si{\USD\per\mega\watt\hour}, base \num{124.5}}}
      & Demand growth        & $+5.5\%$  & $-5.0\%$  & \num{13.09} \\
      & Generation lead time & $-4.7\%$  & $-0.3\%$  & \num{5.52}  \\
      & Transmission lead time & $-1.1\%$ & $+1.0\%$ & \num{2.63}  \\
      & Fuel price           & $+0.4\%$  & $+0.4\%$  & \num{0.01}  \\
    \bottomrule
  \end{tabular}
\end{table}

The result is confirmed when all four factors move together
(Table~\ref{tab:spearman}): demand growth reaches $0.985$ for transmission length and
$0.973$ for installed capacity, close to the maximum a rank correlation can attain,
while no other factor exceeds $0.27$ in absolute value. The two methods agree, which is
the check that matters --- a one-at-a-time result surviving only in the absence of
interaction would not support the conclusion.

\begin{table}[tb]
  \centering
  \caption{Spearman rank correlation between each factor and each indicator over the
  120 Monte Carlo runs, in which all four factors vary simultaneously.}
  \label{tab:spearman}
  \small
  \begin{tabular}{@{}lrrrr@{}}
    \toprule
    Factor & Transmission & Transfer & Capacity & Price \\
           & [\si{\kilo\meter}] & [\si{\giga\watt}] & [\si{\giga\watt}] & [\si{\USD\per\mega\watt\hour}] \\
    \midrule
    Demand growth          & $\mathbf{0.985}$ & $\mathbf{0.835}$ & $\mathbf{0.973}$ & $\mathbf{-0.752}$ \\
    Transmission lead time & $0.112$  & $-0.142$ & $0.121$  & $-0.081$ \\
    Generation lead time   & $0.044$  & $0.081$  & $0.053$  & $-0.264$ \\
    Fuel price             & $-0.008$ & $-0.013$ & $-0.014$ & $-0.014$ \\
    \bottomrule
  \end{tabular}
\end{table}

\subsection{How wide the range of outcomes remains}
\label{sec:results:bands}

Figure~\ref{fig:bands} shows the P10--P90 envelope of each indicator. The envelopes are
narrow until the early 2030s and widen steadily thereafter --- the expected signature of
a system whose divergence is driven by a growth rate rather than a level.

\begin{figure*}[tb]
  \centering
  \includegraphics[width=\textwidth]{01_distribution_P10_P90_timeseries.png}
  \caption{P10--P90 envelope of the four indicators over 120 Monte Carlo runs, with
  the ensemble median and the central case. The band contains the central 80\,\% of
  outcomes.}
  \label{fig:bands}
\end{figure*}

By 2050 the ranges are substantial (Table~\ref{tab:bands}): installed capacity spans
\SIrange{52.3}{82.5}{\giga\watt} between the tenth and ninetieth percentile, a factor
of 1.6, and required transmission \SIrange{39858}{57069}{\kilo\meter}. The price range
is comparatively narrow, \SIrange{116.8}{129.4}{\USD\per\mega\watt\hour}. That
asymmetry is itself informative: what these uncertainties move is \emph{how much
infrastructure the country has to build}, far more than \emph{what electricity ends up
costing}.

\begin{table}[tb]
  \centering
  \caption{Indicator values at January 2050 over the 120 Monte Carlo runs.}
  \label{tab:bands}
  \small
  \begin{tabular}{@{}lrrrrr@{}}
    \toprule
    Indicator & Central & P10 & P50 & P90 & P90$-$P10 \\
    \midrule
    Transmission [\si{\kilo\meter}]            & \num{46341} & \num{39858} & \num{47788} & \num{57069} & \num{17211} \\
    Transfer capability [\si{\giga\watt}]      & \num{4.56}  & \num{4.23}  & \num{4.70}  & \num{5.29}  & \num{1.06}  \\
    Installed capacity [\si{\giga\watt}]       & \num{63.85} & \num{52.33} & \num{66.82} & \num{82.49} & \num{30.16} \\
    Average price [\si{\USD\per\mega\watt\hour}] & \num{124.5} & \num{116.8} & \num{122.8} & \num{129.4} & \num{12.59} \\
    \bottomrule
  \end{tabular}
\end{table}

The central case sits close to the median on transmission and capacity but at the
upper end of the price distribution. That is a consequence of the price
non-monotonicity identified in Table~\ref{tab:tornado}: because both an increase
and a decrease in several factors lower the price relative to the reference, the
reference is near a local maximum.

Supplementary Material~S4.2 reports the convergence of these
percentiles as runs are added. P10, P50 and P90 are stable from roughly 75 runs
onwards for every indicator, which is the evidence that 120 samples are sufficient
rather than an assumption that they are.

\subsection{Central, best and worst case}
\label{sec:results:extremes}

Table~\ref{tab:extremes} reports the three trajectories selected as described in
Section~\ref{sec:method:scenarios}. The best case on delivered price is run
\texttt{mc\_009}, in which demand grows 17.3\,\% faster than the reference; the
worst is \texttt{mc\_058}, in which it grows 16.5\,\% slower.

\begin{table}[tb]
  \centering
  \caption{Central, best and worst case. Best and worst are the Monte Carlo runs
  with the lowest and the highest country average price over the final twelve
  months. Percentages are relative to the central case.}
  \label{tab:extremes}
  \small
  \begin{tabular}{@{}lrrr@{}}
    \toprule
    & Central & Best (\texttt{mc\_009}) & Worst (\texttt{mc\_058}) \\
    \midrule
    \multicolumn{4}{@{}l}{\emph{Factor values}} \\
    Demand growth           & $0.0\%$ & $+17.3\%$ & $-16.5\%$ \\
    Generation lead time    & $0.0\%$ & $+9.8\%$  & $-5.4\%$  \\
    Transmission lead time  & $0.0\%$ & $+2.6\%$  & $+8.2\%$  \\
    Fuel price              & $0.0\%$ & $-2.5\%$  & $-1.1\%$  \\
    \midrule
    \multicolumn{4}{@{}l}{\emph{Indicators at January 2050}} \\
    Transmission [\si{\kilo\meter}]       & \num{46341} & \num{58562} $(+26.4\%)$ & \num{38927} $(-16.0\%)$ \\
    Transfer capability [\si{\giga\watt}] & \num{4.56}  & \num{5.03} $(+10.3\%)$  & \num{4.34} $(-4.8\%)$ \\
    Installed capacity [\si{\giga\watt}]  & \num{63.85} & \num{86.40} $(+35.3\%)$ & \num{49.18} $(-23.0\%)$ \\
    Average price [\si{\USD\per\mega\watt\hour}] & \num{124.5} & \num{114.9} $(-7.6\%)$ & \num{135.6} $(+9.0\%)$ \\
    \bottomrule
  \end{tabular}
\end{table}

This is the central empirical finding, and it runs against the intuition that higher
demand means a more stressed and therefore more expensive system. The mechanism is
visible in the technology mix (Figures~\ref{fig:captech} and~\ref{fig:gentech}): in
the best case onshore wind reaches \SI{52.05}{\giga\watt} against
\SI{34.83}{\giga\watt} centrally and \SI{25.47}{\giga\watt} in the worst, and solar PV
\SI{11.51}{\giga\watt} against \SI{8.33} and \SI{5.24}{\giga\watt}. Faster demand
growth sustains a price signal long enough for low-cost renewable capacity to clear the
profitability threshold earlier and in greater volume, and once in place that capacity
depresses the marginal price. Slow demand growth suppresses the investment signal and
leaves the system leaning on thermal plant --- \SI{10.2}{\tera\watt\hour\per\year} in
the worst case against \SI{5.6}{\tera\watt\hour\per\year} in the best --- ending with a
higher average price despite a smaller system.

\begin{figure*}[tb]
  \centering
  \includegraphics[width=\textwidth]{05_installed_capacity_by_technology.png}
  \caption{Installed capacity by technology for the central, best and worst case.
  The dashed line is national peak demand. Fourteen model technologies are
  aggregated into eight series; the mapping is given in Supplementary
  Material~S5.3.}
  \label{fig:captech}
\end{figure*}

\begin{figure*}[tb]
  \centering
  \includegraphics[width=\textwidth]{06_generation_by_technology.png}
  \caption{Annual generation by technology for the central, best and worst case.}
  \label{fig:gentech}
\end{figure*}

The corollary is uncomfortable for a planner. The least-cost trajectory is also the
most infrastructure-intensive: the best case requires \SI{58562}{\kilo\meter} of
network, 26.4\,\% more than the central case and 50\,\% more than the worst case,
together with 10.3\,\% more inter-zonal transfer capability. The cheapest electricity
in this system is not obtained by building less; it is obtained by building more,
faster. Figure~\ref{fig:price} shows the resulting price paths and
Figure~\ref{fig:transmission} the transmission requirement.

\begin{figure*}[tb]
  \centering
  \includegraphics[width=\textwidth]{04_price_best_worst.png}
  \caption{Country average electricity price for the central, best and worst case.}
  \label{fig:price}
\end{figure*}

\begin{figure*}[tb]
  \centering
  \includegraphics[width=\textwidth]{08_transmission_km.png}
  \caption{Transmission network by category --- demand connection, generation
  connection and interconnection --- for the central, best and worst case.}
  \label{fig:transmission}
\end{figure*}

A second observation concerns the delay factors. The best case has \emph{longer}
generation lead times than the central case ($+9.8$\,\%) and the worst case has
shorter ones ($-5.4$\,\%). This is not evidence that delays are beneficial; it is a
consequence of their weakness relative to demand. Within the sampled ranges, the
demand draw determines the outcome and the delay draw is close to noise, so the
extreme runs on price carry essentially arbitrary delay values. Reporting this
explicitly is more honest than selecting extreme runs that happen to tell a tidier
story, and it is itself the clearest statement of the ranking: in this system, over
these ranges, lead times do not decide the outcome.

\subsection{Where the uncertainty acts}
\label{sec:results:where}

Decomposing transmission by category (Figure~\ref{fig:transmission}) locates the
effect. Network built for demand connection scales directly with peak demand and
inherits the demand factor's variance in full; network built to connect generation
follows the investment decision and inherits it indirectly. Inter-zonal corridor
capability is the least sensitive of the three and the only one where transmission
lead time ranks second rather than fourth (Table~\ref{tab:tornado}): a longer pipeline
defers reinforcements, which the series topology converts into persistent congestion
rather than a different final capacity. That is why transmission lead time barely
registers in the aggregate indicators while remaining operationally important --- the
model builds much the same network in the end, and what the lead time changes is
\emph{when} it is available.
